% !TEX program = xelatex
% 库的相关工作 —— 前提选择与引理库演化综述
\documentclass[UTF8,a4paper,11pt,fontset=windows]{ctexart}

\usepackage{geometry}
\geometry{left=2.3cm,right=2.3cm,top=2.5cm,bottom=2.3cm}
\usepackage{amsmath,amssymb}
\usepackage{booktabs}
\usepackage{tabularx}
\usepackage{array}
\usepackage{enumitem}
\usepackage{xcolor}
\usepackage{url}
\usepackage{adjustbox}
\usepackage[hidelinks]{hyperref}
\usepackage{fancyhdr}
\usepackage{tcolorbox}
\tcbuselibrary{skins,breakable}

\definecolor{sgblue}{RGB}{31,86,153}
\definecolor{sggreen}{RGB}{38,120,86}
\definecolor{sggray}{RGB}{90,96,106}
\definecolor{sgred}{RGB}{150,50,50}

\xeCJKDeclareCharClass{CJK}{"2460 -> "24FF}

\ctexset{
  section/format = \Large\bfseries\color{sgblue},
  subsection/format = \large\bfseries\color{sgblue!85!black},
  subsubsection/format = \normalsize\bfseries\color{sggray},
}

\pagestyle{fancy}
\fancyhf{}
\fancyhead[L]{\small\color{sggray}库的相关工作}
\fancyhead[R]{\small\color{sggray}\leftmark}
\fancyfoot[C]{\small\thepage}
\renewcommand{\headrulewidth}{0.4pt}

\newcolumntype{Y}{>{\raggedright\arraybackslash}X}
\newcommand{\lp}[1]{\path{#1}}

\newtcolorbox{keybox}[1][]{
  colback=blue!4, colframe=sgblue!70, boxrule=0.6pt, arc=2pt,
  left=6pt, right=6pt, top=5pt, bottom=5pt, breakable, #1}
\newtcolorbox{warnbox}[1][]{
  colback=red!3, colframe=sgred!60, boxrule=0.6pt, arc=2pt,
  left=6pt, right=6pt, top=5pt, bottom=5pt, breakable, #1}

\newcommand{\paper}[1]{\par\medskip\noindent\textbf{\color{sgblue!80!black}#1}\par\nobreak\vspace{1pt}}
\newcommand{\field}[1]{\textit{#1}：}

\title{\bfseries 引理库相关工作综述\\[2mm]
       \large 前提选择（Premise Selection）与引理库演化（Library Learning）}
\author{}
\date{}

\begin{document}
\maketitle
\thispagestyle{fancy}

\begin{keybox}
\textbf{文档说明}：本文件是为项目的"引理库的建立与管理"模块所做的文献调研，
覆盖\textbf{前提选择}与\textbf{引理库演化}两条技术线。
检索方式为 arXiv API 关键词检索（主题词包括 premise selection、retrieval + theorem proving、
lemma library、submodular 等），并对关键论文下载全文精读；
本次共下载并入藏 22 篇，清单见附录 A。
本文件与《SG-Lean 思路汇报》配套阅读：后者讲我们打算怎么做，本文件讲别人已经做到了哪一步。
\end{keybox}

\vspace{2mm}
\tableofcontents
\vspace{4mm}

%======================================================================
\section{综述范围与方法}

\subsection{我们要回答的问题}

项目在"引理库"这条线上原本的设想可以拆成四个动作：
（1）\textbf{建}：从一个课程语料出发，产出候选引理；
（2）\textbf{选}：在库容有限的前提下决定留下哪些；
（3）\textbf{管}：排序、去重、淘汰、版本；
（4）\textbf{用}：在线解题时把库里的引理喂给证明器。

本综述要回答的是：这四个动作，学术上分别被谁做过了，做到什么程度，用的什么判据，报了什么数字。
只有把这条基线摸清，才能判断我们剩下的可主张范围是什么。

\subsection{一个关键澄清：这是两条线，不是一条}

检索过程中很快会发现，"引理库"在文献里被拆成两个独立的研究方向，各自都有十年以上的积累：

\begin{table}[htbp]
\centering
\small
\begin{tabularx}{\textwidth}{@{}p{2.6cm}p{5.2cm}Y@{}}
\toprule
\textbf{技术线} & \textbf{回答的问题} & \textbf{典型输出} \\
\midrule
\textbf{线 A：前提选择} & 给定一个证明状态或命题，从一个大库里挑出对证明有用的前提 &
一个排序好的前提列表（每次查询返回前 $k$ 条） \\
\midrule
\textbf{线 B：引理库演化} & 主动发现、抽象、压缩、淘汰引理，形成一个小而精的库 &
一个可持续增长的引理库本身 \\
\bottomrule
\end{tabularx}
\end{table}

两条线的交汇点在"用"：线 B 产出库，线 A 决定库里哪几条进上下文。
我们的设想同时压在这两条线上，所以两条都要看。

\subsection{一个被反复引用的痛点}

这条赛道上有一个反复出现的经验结论，值得单独记住，因为它是我们项目的立论基础之一：

\begin{keybox}
从证明过程中自动抽取出来的引理，\textbf{大多问题特异、不可复用}。
LEGO-Prover 在 2023 年明确写出这一点，并以此作为设计动机；
DreamProver 在 2026 年仍以"现有方法要么依赖固定库，要么合成的中间引理高度特异、缺乏通用性"作为出发点。
\end{keybox}

也就是说，"抽取的引理不可复用"是八年前就被看见、至今仍被当作开放问题的事。
这既说明我们的问题意识是对的，也说明这条路上已经站着很多人。

%======================================================================
\section{前提选择（线 A）}

\subsection{奠基：把前提选择建成二分类问题}

\paper{DeepMath, arXiv:1606.04442, 2016（Google，含 Urban）}
\field{做法}把前提选择形式化为\textbf{二分类}：将待证命题与每个候选前提分别嵌入，
用卷积、循环或混合网络编码，再拼接送入一层 logit 判断相关性。
\field{意义}这是"\textbf{用神经网络做前提选择}"的开创性工作，
确立了此后十余年的基本范式：\textbf{嵌入 + 相关性判别}。
\field{对我们的意义}我们的"库"如果只是一个可检索的引理集合，
那么它的检索方式必须与这套范式比较，而我们的起点是零（目前连检索都没有）。

\paper{FormulaNet 及图网络系列, 2017 至 2022}
\field{做法}形式化语句天然是树或图。此后一条支流把语句解析成树或图，
用树结构网络或\textbf{图神经网络}编码：FormulaNet 用保留边序的图嵌入；
Olšák 等人在一阶逻辑公式上构造能捕捉多种逻辑不变性的图网络框架；
Paliwal 等人系统比较了 HOL Light 中公式的多种图表示；后续工作引入\textbf{图对比学习}；
Ferreira 与 Freitas 以及 Bauer 等人则在大规模语料上构造\textbf{依赖图}，
把定理与前提作为节点、依赖关系作为边，用图网络预测节点之间的连边来完成前提选择。
\field{对我们的意义}"依赖图"这个结构与我们想用的"引理被谁用到"的统计高度相关，
需要明确我们与它的区别：他们用图结构做\textbf{预测}，我们想做的是用\textbf{实测覆盖增益}做筛选。

\subsection{神经引导的符号证明器}

\paper{The Isabelle ENIGMA, arXiv:2205.01981, 2022}
\field{做法}把学习方法与定理证明器紧密结合：为 Isabelle 的 Sledgehammer 问题训练
\textbf{定制的 ENIGMA 引导}、\textbf{定制的神经前提选择}，并为 E 证明器设计针对性策略。
训练数据是从 Isabelle 中抽取的数十万条无类型与有类型一阶问题，并\textbf{迭代多轮}训练。
\field{结果}最终最好的单策略 ENIGMA 与前提选择系统，在 15 秒内
\textbf{比此前最好的 E 提升 25.3\%}，同时超过所有已有的 ATP 与 SMT 系统。
\field{对我们的意义}这是"学习式前提选择确实能显著提升证明能力"的早期强证据，
也说明这类增益必须在\textbf{统一时间预算}下比较才有意义（他们明确写了 15 秒）。

\subsection{对比学习与 Transformer}

\paper{Magnushammer, arXiv:2303.04488, 2023}
\field{做法}用 Transformer 与\textbf{对比学习}做前提检索，明确以"避免符号方法与工程开销"为卖点。
\field{结果}PISA 基准 \textbf{59.5\% 对 Sledgehammer 的 38.3\%}；
miniF2F 基准 \textbf{34.0\% 对 20.9\%}；
再与语言模型证明器结合，把 PISA 上的最好证明成功率从 57.0\% 提到 \textbf{71.0\%}，
而参数量少四倍。此外发布了当时最大规模的前提选择数据集（证明状态与相关前提的文本对）。
\field{对我们的意义}"检索做得好，末端证明成功率就能被推上去"在这里已有量化证据；
也说明\textbf{检索器本身的训练数据规模}是这类方法的关键资源。

\subsection{检索增强证明的范式确立}

\paper{LeanDojo 与 ReProver, arXiv:2306.15626, 2023}
\field{做法}提出检索增强的 Lean 证明器 ReProver。
关键设计是用程序分析识别\textbf{可访问前提}并构造\textbf{困难负样本}，
使检索器的训练更有效（基于稠密段落检索 DPR）。
同时构造了包含 98{,}734 条定理的基准。
\field{结果}检索显著优于无检索版本。
\field{对我们的意义}"可访问前提 + 困难负样本"是训练检索器的标准配方；
如果我们要自己做检索，这是必须复用的工程经验。

\subsection{Lean 生态的检索工具爆发（2023 至 2026）}

这是与我们最直接相关的一节。近三年 Lean 侧的检索工具密集出现，功能已经相当完整。

\paper{Machine-Learned Premise Selection for Lean, arXiv:2304.00994, 2023}
\field{做法}在 \textbf{Lean 内部}实现在线训练的\textbf{随机森林}，数据从 Mathlib 抽取，
以 \lp{suggest_premises} tactic 的形式在编辑器里供用户调用。
设计原则写得很清楚：与证明助手紧耦合、易安装、轻量且快。
\field{对我们的意义}"轻量、可嵌入、在证明助手内直接可用"是一条已被验证的产品化路线。

\paper{LeanExplore, arXiv:2506.11085, 2025}
\field{做法}Lean 声明检索引擎，\textbf{混合排序}：多源语义嵌入（形式代码、docstring、
AI 生成的非形式化翻译、声明标题）\textbf{加 BM25+ 词法相关性}\textbf{加基于 PageRank 的重要度分数}
（用声明的相互连接程度反映重要性）。覆盖 Batteries、Init、Lean、Mathlib、PhysLean、Std。
提供网站、Python API 与 MCP，可自托管。
\field{对我们的意义}\textbf{这一条直接对应我们设想中的"给常用的引理更高的地位"}。
他们用的是图上的 PageRank，而不是实测使用频率。这是一个可以直接对照的判据差异。

\paper{Lean Finder, arXiv:2510.15940, 2025}
\field{做法}面向\textbf{数学家意图}的语义搜索：先分析并聚类公开 Lean 讨论，
再用合成查询微调文本嵌入，并用多种反馈信号对齐用户偏好。
\field{结果}在真实查询、非形式化陈述与证明状态三类评测上，
相对此前检索引擎与 GPT-4o \textbf{相对提升超过 30\%}。
\field{对我们的意义}检索质量的下一个瓶颈是"查询与库表示之间的意图错配"，
这条对我们设计提示词与查询构造有直接参考价值。

\paper{LeanPremise 与 LeanHammer, arXiv:2506.07477, 2025}
\field{做法}LeanPremise 是首个\textbf{面向 hammer 使用场景、针对依赖类型论训练}的前提选择器。
关键特性是\textbf{动态适应用户上下文}：运行时从用户环境动态抽取签名，
从而推荐训练数据之外的库以及用户本地定义的引理。
\field{结果}相比已有前提选择器，使 LeanHammer \textbf{多解 21\% 的目标}；
并在 miniCTX-v2 上评测了泛化能力。
\field{对我们的意义}\textbf{"动态适应用户上下文"几乎就是我们要做的事}：
我们的库本质上就是"训练数据之外的、动态积累的引理集合"。这一点必须正面处理。

\paper{LeanSearch v2, arXiv:2605.13137, 2026}
\field{问题定义}论文明确指出既有工作只覆盖了两个相邻问题：
语义检索引擎找的是\textbf{单条}匹配声明的声明；
前提选择系统预测的是\textbf{单步 tactic} 需要的前提。
两者都不能给出"一整条定理所需的前提集合"。作者把这件事命名为\textbf{全局前提检索}。
\field{做法}两模式。标准模式：构建\textbf{层级化非形式化}的 Mathlib 语料，
用嵌入加重排的两阶段流水线，且\textbf{不做领域微调}。
推理模式：以标准模式为检索底座，做迭代的\textbf{草图、检索、反思}循环，
为目标定理生成证明草图，再为草图的每一步检索候选、过滤并评判，迭代修订草图直到得到一致的支撑引理集合。
\field{结果}标准模式 nDCG@10 达 \textbf{0.62}（次优系统 0.53）。
在 69 条研究级 Mathlib 定理查询的基准上，推理模式在 10 个候选内恢复
\textbf{46.1\%} 的真实前提组（对比：推理式检索器 38.0\%，前提选择基线 9.3\%）。
\field{下游验证（最重要的一条）}在固定的证明循环里替换检索器：
LeanSearch v2 取得最高证明成功率 \textbf{20\%}，次优系统 16\%，\textbf{无检索只有 4\%}。
\field{对我们的意义}这是全篇最值得我们记住的数字。
它同时说明两件事：（1）\textbf{检索质量确实直接传导为证明成功率}，杠杆很大；
（2）"按需检索出相关引理集合"这件事已经被做到相当深，不能作为我们的创新点。

\paper{Combining Textual and Structural Information, arXiv:2510.23637, 2025}
\field{做法}图增强的前提选择：用稠密文本嵌入结合\textbf{异构图}上的图神经网络，
图中同时包含"状态与前提"和"前提与前提"两类关系。
\field{结果}在 LeanDojo 基准上，相对 ReProver 的语言基线，标准检索指标\textbf{提升超过 25\%}。
\field{对我们的意义}再次说明"N 条引理之间的关系（依赖、共现）"是有效信号，
而这恰好也是"库的结构化组织"能提供的东西。

\paper{Semantic Search over 9 Million Mathematical Theorems, arXiv:2602.05216, 2026}
\field{做法}在规模上做文章：把 arXiv 等八个来源的 \textbf{920 万}条定理陈述统一成语料，
用"简短自然语言描述"作为检索表示，系统分析表示上下文、语言模型选择、嵌入模型与提示策略的影响。
\field{结果}在职业数学家撰写的查询集上，定理级与论文级检索都大幅优于已有基线。
\field{对我们的意义}说明"为每条陈述生成一句自然语言描述再检索"是一个稳定有效的表示，
我们若要给库中的引理建索引，这是可直接借鉴的做法。

\paper{CSLibPremiseBench, arXiv:2605.14549, 2026（审计型负面结果）}
\field{做法}针对 Lean 4 的计算机科学库 CSLib 构建可复现的前提检索基准：
固定 CSLib v4.29.0，任务集含 801 条可标注的定理与引理任务、1875 条候选声明。
对比 BM25、符号与命名重叠、命名空间与模块启发式、PageRank 式模块先验、
固定混合方案，以及作者自提出的结构引导图词法重排器 CSG-Rerank。
\field{标签可靠性审计}他们做的这件事很值得学：
标签是"源码可见的证明引用代理"而非 elaborate 后的依赖轨迹；
严格审计保留 2666/2845（93.7\%）的代理链接，
但用 300 条任务做 Lean 环境表达式探针时，只在 958 条代理链接中观察到 462 条（48.2\%）。
\field{结果}结构引导的重排只在严格策略下带来微弱的早期排名提升，
\textbf{无法可靠超过 BM25 加符号重叠}，Recall@10 也没有改善。
\field{对我们的意义}\textbf{两条警示}：一是\textbf{朴素的结构化启发式（命名空间、模块、图先验）不可靠}，
别指望"领域划分 + 命名规则"就能提升检索；二是\textbf{标签本身可能非常不可靠}，
任何关于"检索命中率"的指标都必须先审计标签。

\subsection{线 A 小结}

\begin{table}[htbp]
\centering
\footnotesize
\begin{tabularx}{\textwidth}{@{}p{3.5cm}p{4.2cm}Y@{}}
\toprule
\textbf{阶段} & \textbf{代表} & \textbf{方法要点} \\
\midrule
二分类奠基 & DeepMath & 嵌入 + 相关性判别 \\
\midrule
结构化表示 & FormulaNet、GNN 系列、依赖图 & 树/图编码，图对比学习，连边预测 \\
\midrule
神经引导符号器 & Isabelle ENIGMA & 学习引导 + 前提选择 + 定制策略，统一时限比较 \\
\midrule
对比学习 & Magnushammer、PACT & Transformer + 对比/自回归目标 \\
\midrule
检索增强证明 & LeanDojo / ReProver & 可访问前提 + 困难负样本 \\
\midrule
Lean 生态工具 & LeanExplore、Lean Finder & 混合排序（语义 + BM25+ + PageRank）、意图对齐 \\
\midrule
端到端 hammer & LeanPremise / LeanHammer & 为依赖类型论训练，动态适应用户上下文 \\
\midrule
全局前提检索 & LeanSearch v2 & 一次给出一组前提，迭代草图与检索，下游增益显著 \\
\midrule
基准与审计 & CSLibPremiseBench & 负面结果：结构启发式不可靠；标签需审计 \\
\bottomrule
\end{tabularx}
\end{table}

\textbf{结论}：线 A 已经成熟到"有公开检索引擎、有 API、有 MCP、有基准、有负面结果"的程度。
凡是"给引理排序、按需检索、动态适配、重要度加权"这类想法，\textbf{都已被占}。

%======================================================================
\section{引理库的演化与管理（线 B）}

这一节与我们的模块最直接相关，按时间顺序展开。

\subsection{有用性驱动的引理挖掘}

\paper{Learning-assisted Theorem Proving with Millions of Lemmas, arXiv:1402.3578, 2014}
\field{做法}大型形式化库里存在数百万条已证语句，但与人写数学一样，
其中只有极小一部分被命名并在后续证明中被复用。
作者\textbf{提出并实现了一套判据}，用来估计 HOL Light 引理对证明后续定理的"有用性"，
并据此在 HOL Light 与 Flyspeck 的\textbf{大规模推理图}中挖掘出最好的引理，
加入可复用的语句池，再配合基于学习的相关性过滤。
\field{结果}这种组合显著增强了在 Flyspeck 这类大型形式化库上对\textbf{新猜想}的自动证明
（作为背景：当时 HOLyHammer 能证明 Flyspeck 中约 47\% 的定理）。
\field{对我们的意义}\textbf{这是全篇最需要警惕的一条}："按有用性给引理打分并筛选入库"
是\textbf{十二年前}就有的想法。我们如果只是把"有用性"换成另一个名字，没有任何新意。

\subsection{成长型技能库}

\paper{LEGO-Prover, arXiv:2310.00656, 2023}
\field{做法}指出现有方法的共同局限是\textbf{假设定理库在整个证明过程中固定不变}。
LEGO-Prover 维护一个\textbf{不断增长的技能库}（已验证的引理作为技能）：
模块化地构造证明，使模型既能\textbf{检索使用}已有技能，又能在证明过程中\textbf{创建新技能}；
这些技能还会被进一步"演化"（通过提示语言模型）以变得更通用、更可复用，
从而在另一个尺度上丰富库。此外，学到的库还能弥合人类证明与形式证明之间的差距，
使缺失步骤更容易被填补。
\field{结果}miniF2F 通过率：valid 从 48.0\% 提到 \textbf{57.0\%}，test 从 45.5\% 提到 \textbf{50.0\%}；
过程中生成\textbf{两万多个}技能并加入库；
消融显示新加入的技能确实有用：成功率从 47.1\% 提升到 \textbf{50.4\%}。
\field{局限（论文自述，也是我们的机会）}从证明器抽取出的引理"大多是问题特异的，因而不可复用"，
所以他们要额外加一个"演化器"来提升通用性。
\field{对我们的意义}这是"库的建立 + 库的使用"的第一篇完整工程实现，
而且它给出了一个\textbf{可直接对照的增益量级}：库的净贡献约 3 个百分点。

\subsection{wake-sleep：把库的演化做成两阶段循环（最接近我们的工作）}

\paper{DreamProver, arXiv:2604.26311, 2026}

这是本次调研中对我们的设想覆盖最完整的一篇，因此单独展开。

\field{问题陈述}现有方法要么依赖\textbf{固定}的引理库（缺乏适应性），
要么合成的中间引理\textbf{高度特异}（缺乏通用性）。

\field{总体框架}用"wake-sleep"程序归纳范式，构成迭代的两阶段循环：

\begin{itemize}[leftmargin=1.5em,itemsep=2pt]
  \item \textbf{wake 阶段}：系统尝试用\textbf{当前库}去证明训练集中的定理，
        同时提出新的引理候选（来自证明过程中的中间定理）。
  \item \textbf{sleep 阶段}：系统从这些经验中学习，把候选\textbf{抽象、精炼、整合}，
        以\textbf{压缩并优化}库。
\end{itemize}

\field{sleep 阶段的第一步：引理抽象}采用基于聚类的演化策略：
先给每条中间定理标注自然语言描述（含子领域、难度，例如"是否只是平凡的等价变换"
或"是否需要更高级的推理"，以及可用于其证明的相关定理）；
再对描述做嵌入并用余弦距离做 \textbf{K-Means 聚类}（簇数用肘部法确定）；
然后对每个簇，把簇内定理及其描述作为上下文交给语言模型，
生成\textbf{更一般、更可复用}的候选引理；
最后用\textbf{结构化相似度}校验候选是否真的适用于该簇：
把候选引理与簇内定理都转成简化的一阶逻辑表达式树，比较树之间的可对齐程度，
最高相似度超过阈值才保留。

\field{sleep 阶段的第二步：库更新（三步）}
\begin{enumerate}[leftmargin=1.5em,itemsep=2pt]
  \item \textbf{遗忘}：追踪每条引理在过去各个 wake 阶段的使用频率，
        采用\textbf{最近最少使用}（LRU）策略，当库达到容量上限时淘汰最少被使用的引理。
        论文明确说明这样做的理由之一是\textbf{避免把整个库作为输入上下文时把重要引理挤掉}。
  \item \textbf{存入非重复项}：用\textbf{树编辑距离}衡量候选与库中已有引理的相似度，
        距离很小（即高度相似）的视为重复并移除，以降低冗余、保持多样性。
  \item \textbf{形式化验证}：对剩下的每条候选引理，用与 wake 阶段相同的直接提示策略尝试证明它。
\end{enumerate}

\field{推理阶段}先用库做直接提示证明；若失败，退化为\textbf{草图与证明}工作流：
先生成证明草图，再逐个证明子目标，两个环节都继续使用库。

\field{关键规模数字}库被维持在\textbf{不到 100 条}引理。
作者称这个遗忘机制"维持了一个简洁的引理库"，从而支持稳定高效的复用。

\field{结果}相对此前的状态最优系统 Hilbert，在不平等、数论、组合三个领域分别提升
20\%、114\%、50\%，\textbf{平均提升 61\%}；
相比 Goedel-Prover-V2，输出 token 分别减少 55\%、56\%、76\%；
相比 Hilbert 减少 42\%、50\%、53\%；
证明长度相比两者平均减少 32\% 与 50\%。
\field{数据设置}库是\textbf{按领域分别构建}的：每个领域从对应的训练子集采 \textbf{100 道题}；
评测用 567NEQ、ChenNEQ、MO-INT（不等式）、PutnamBench 与 ProverBench 的数论子集、
CombiBench 子集（组合）。
\field{泛化范围}论文用单独的研究问题考察"低资源领域的泛化"，
但表述始终是"\textbf{相关领域的未见定理}"，并未声称跨远域泛化。

\begin{warnbox}
\textbf{对我们的三点直接含义}：
\begin{enumerate}[leftmargin=1.5em,itemsep=2pt,topsep=2pt]
  \item 我们设想的"库的建立、演化、淘汰、去重、验证"这条链，
        \textbf{已经被完整实现过一遍}，而且每一环都有具体判据：
        聚类与抽象用于选择，LRU 使用频率用于排序与淘汰，
        树编辑距离用于去重，形式化验证用于准入。
  \item 他们的\textbf{选择判据是启发式}：语义聚类与结构相似度阈值，
        \textbf{不是实测的工作负载覆盖增益}，也没有近似保证。差异点正在这里。
  \item 他们被迫把库压到 100 条以内，原因是\textbf{整库要作为上下文输入}。
        他们选择压缩库来适配上下文；我们设想的是\textbf{用检索从库里取前若干条}。
        这是两种不同的解法，而后者正是线 A 已经证明有效的那一条。
\end{enumerate}
\end{warnbox}

\subsection{folklore 引理的生产线}

\paper{MathlibLemma, arXiv:2602.02561, 2026}
\field{做法}面向一个具体而实用的问题：Mathlib 里缺少大量\textbf{民间引理}
（数学家习以为常、但形式库中不一定有的中间事实），这限制了 Lean 作为日常工具的可用性。
作者提出模块化流水线，自动完成这类引理的\textbf{发现、形式化与证明}。
\field{结果}产出一个经\textbf{验证的引理库}，包含 \textbf{1506 条}通过"证明旁路筛查"
（proof-bypass screen）的 Lean 检查通过的证明；
其中一个小的精选子集\textbf{已被合并进 Mathlib}，作为"输出能达到专家库标准"的外部证据。
此外构造了 MathlibLemma 基准：\textbf{4028 条}非平凡、通过类型检查的语句，
覆盖广泛数学领域；对未被证明的残差做种子分层抽样审计，\textbf{78\%} 在数学上成立
（说明这些未解实例大多是有效命题而非幻觉）。
在该基准上，所有模型合计证明 37\%，而\textbf{单个模型最多约 19\%}。
\field{对我们的意义}这是"自造引理入库"的当代最强形态，而且它证明了这类产出
\textbf{能达到被人工库接纳的质量}。同时它给出了一个有用的基准与一个可对照的口径：
单模型上限约 19\%，集合约 37\%，说明这类任务仍有很大空间。

\subsection{把新验证的证明索引进检索记忆}

\paper{OProver, arXiv:2605.17283, 2026}
\field{做法}指出智能体式证明很少被纳入训练，通常只在推理时使用。
OProver 把二者统一：失败的证明尝试用"检索到的、经编译器验证的证明"
与 Lean 编译器反馈\textbf{迭代修正}；每一轮迭代都执行智能体式证明，
把\textbf{新验证的证明索引进} proof 集合与\textbf{检索记忆}，
把修复轨迹用作监督微调数据，把未解的难题用于强化学习。
\field{规模与结果}其 OProofs 集合含 177 万条 Lean 语句、686 万条编译器验证的证明，
以及带有检索上下文、失败尝试、反馈与修复过程的序列化轨迹。
在五个基准上取得最好或次好的成绩，其中 MiniF2F 为 93.3\%（Pass@32）。
\field{对我们的意义}它把"新证明入库"和"检索记忆"做成了训练闭环的一部分，
说明\textbf{库与检索记忆的持续增长是当前主流系统的标准组件}。

\subsection{分解式证明：库的另一种用法}

\paper{Hilbert, arXiv:2509.22819, 2025}
\field{做法}编排四个组件：擅长度推理的非形式化语言模型、
为 Lean 4 优化的专用证明模型、形式验证器、\textbf{语义定理检索器}。
当证明模型无法解决时，用\textbf{递归分解}把问题拆成子目标，
交给证明模型或推理模型解决，并利用验证器反馈精修错误证明。
\field{结果}miniF2F 达 \textbf{99.2\%}（比当时公开最好高 6.6 个百分点）；
PutnamBench 解出 462/660（\textbf{70.0\%}）。
\field{对我们的意义}\textbf{这个数字很重要}：miniF2F 已经被做到 99.2\%，
意味着\textbf{该基准几乎没有增益空间}。我们原定的"在 miniF2F 上测增益"的方案必须修改。

\subsection{领域专用的可复用证明库}

\paper{CircuitProver, arXiv:2607.27259, 2026}
\field{做法}硬件形式验证场景：把参数化硬件设计与自然语言规格转成可执行的 Lean 4 模型，
再通过 Lean 反馈迭代构造机器检查的证明；把\textbf{证明轨迹与已验证定理蒸馏成可复用的库}，
库中的证明策略指导后续智能体的推理，已验证引理可被复用；支持证明积累与参数化验证。
\field{对我们的意义}这是"库"在\textbf{垂直领域}里的应用形态，
说明按领域建库在实践中是成立的（与 DreamProver 的按领域建库一致）。

\subsection{线 B 小结：库的生命周期对照表}

\begin{table}[htbp]
\centering
\footnotesize
\begin{tabularx}{\textwidth}{@{}p{2.6cm}p{4.0cm}Y@{}}
\toprule
\textbf{库生命周期的环节} & \textbf{谁做过} & \textbf{用的判据} \\
\midrule
建：产出候选引理 & LEGO-Prover、DreamProver、MathlibLemma、OProver &
证明过程中的中间定理；针对民间引理的定向挖掘；智能体式证明 \\
\midrule
选：决定留下哪些 & DreamProver & 语义聚类 + 抽象 + 结构化相似度阈值（\textbf{启发式}） \\
\midrule
准入：能不能进 & DreamProver、MathlibLemma & 形式化验证 + 证明旁路筛查 \\
\midrule
去重 & LEGO-Prover、DreamProver & 字符串相似度；树编辑距离 \\
\midrule
排序与淘汰 & DreamProver & \textbf{LRU}（跨轮次使用频率） \\
\midrule
容量 & DreamProver & 硬上限，约 100 条，受提示词上下文约束 \\
\midrule
用：进上下文 & DreamProver & \textbf{整库注入}（所以必须压小） \\
\midrule
用：按需检索 & Lean 生态全体（见第三节） & 语义嵌入 + 词法 + 图重要度等混合排序 \\
\bottomrule
\end{tabularx}
\end{table}

\textbf{两个空白可以直接看出来}：
（1）\textbf{没有任何工作用"实测的工作负载覆盖增益"来做选与汰}，全是启发式；
（2）\textbf{库的容量上限是被"整库注入"这个用法逼出来的}，
如果改用检索，容量约束的性质就完全不同（从"能塞进上下文"变成"检索精度"）。

%======================================================================
\section{两线的交汇与横向对照}

\subsection{按"库的使用方式"分类}

\begin{table}[htbp]
\centering
\small
\begin{tabularx}{\textwidth}{@{}p{4.6cm}Y@{}}
\toprule
\textbf{用法} & \textbf{代表与代价} \\
\midrule
\textbf{整库注入提示词} & DreamProver。代价：库必须保持很小（他们在 100 条以内） \\
\midrule
\textbf{逐步检索（tactic 级）} & ReProver、LeanPremise、Magnushammer。代价：需要训练检索器 \\
\midrule
\textbf{全局检索（一次给一组）} & LeanSearch v2 的推理模式。代价：迭代多次查询，调用成本上升 \\
\midrule
\textbf{把新证明索引进记忆} & OProver。代价：需要持续维护索引与训练闭环 \\
\bottomrule
\end{tabularx}
\end{table}

\subsection{评测协议对照}

\begin{table}[htbp]
\centering
\footnotesize
\begin{tabularx}{\textwidth}{@{}p{3.6cm}p{4.0cm}Y@{}}
\toprule
\textbf{工作} & \textbf{基准} & \textbf{关键数字} \\
\midrule
Magnushammer & PISA、miniF2F & 59.5\% 对 38.3\%；34.0\% 对 20.9\% \\
\midrule
Isabelle ENIGMA & Sledgehammer 问题集 & 15 秒内相对提升 25.3\% \\
\midrule
LEGO-Prover & miniF2F & valid 48.0 到 57.0；test 45.5 到 50.0；消融 +3.3 点 \\
\midrule
DreamProver & 不等式、数论、组合 & 相对 Hilbert 平均 +61\%；token 减少最多 76\% \\
\midrule
MathlibLemma & 自建 4028 条基准 & 全模型合计 37\%，单模型最多约 19\% \\
\midrule
Hilbert & miniF2F、PutnamBench & 99.2\%；70.0\% \\
\midrule
OProver & 五个基准 & MiniF2F 93.3\%、PutnamBench 11.3\% \\
\midrule
LeanSearch v2 & 69 条研究级查询 & nDCG@10 0.62；下游证明 20\% 对 16\% 对无检索 4\% \\
\midrule
CSLibPremiseBench & CSLib & 结构重排无法可靠超过 BM25 加符号重叠 \\
\bottomrule
\end{tabularx}
\end{table}

%======================================================================
\section{从这条赛道得到的六条硬结论}

\begin{enumerate}[leftmargin=1.6em,itemsep=4pt]
  \item \textbf{检索是已被验证的大杠杆，而我们还没有。}
        LeanSearch v2 的下游对照给出了 20\% 对 4\%（无检索）的差距。
        我们的流水线目前连检索层都没有，这是最应该先补的一环。
        但正因为它成熟，\textbf{不能作为创新点}。
  \item \textbf{"库的建、选、管、用"每一环都有代表作。}
        DreamProver 是最完整的一次实现，覆盖了聚类抽象、LRU 淘汰、树编辑距离去重、
        形式化验证准入与整库注入。我们原设想的清单基本被逐项覆盖。
  \item \textbf{选择标准全部是启发式。}
        使用频率（LRU）、语义聚类、结构相似度阈值、字符串相似度。
        \textbf{没有工作用"实测的工作负载覆盖增益"并给出近似保证}。
        这可能是一条真实的空位。
  \item \textbf{整库注入是容量瓶颈的根源。}
        DreamProver 明确因为"整库要进上下文"而把库限制在 100 条以内。
        这意味着"库能长多大"被提示词预算而非可用引理数量决定。
        改用检索替代压缩，是一个结构性差异。
  \item \textbf{miniF2F 已饱和，不能作为我们的主基准。}
        Hilbert 99.2\%、OProver 93.3\%，天花板效应会淹没任何真实增益。
        需要改用有提升空间的基准（PutnamBench、CombiBench、ProofNet、ProverBench、
        FATE-M、567NEQ 等）。
  \item \textbf{负面结果与标签审计必须纳入设计。}
        CSLibPremiseBench 说明结构启发式（命名空间、模块、图先验）不可靠，
        且"源码可见的引用代理"标签与真实依赖只有约一半吻合。
        我们若报"库命中率"或"引用率"，必须先做标签可靠性审计。
\end{enumerate}

%======================================================================
\section{我们的定位}

\subsection{已被占用的部分（不要声称）}

\begin{table}[htbp]
\centering
\small
\begin{tabularx}{\textwidth}{@{}p{6.2cm}Y@{}}
\toprule
\textbf{想法} & \textbf{已被谁做} \\
\midrule
从大库里按相关度检索前提 & ReProver、Magnushammer、LeanPremise、LeanExplore、LeanSearch v2 \\
\midrule
给引理按重要度排序 & LeanExplore 用 PageRank；DreamProver 用使用频率 LRU \\
\midrule
按需检索而非整库注入 & Lean 生态的检索工具普遍如此 \\
\midrule
自造引理并入库 & LEGO-Prover、DreamProver、MathlibLemma \\
\midrule
库的去重 & LEGO-Prover 用字符串相似度；DreamProver 用树编辑距离 \\
\midrule
库的淘汰 & DreamProver 的 LRU 遗忘 \\
\midrule
库的压缩与抽象 & DreamProver 的 sleep 阶段（比我们设想的更深） \\
\midrule
把新验证的证明索引进可用记忆 & OProver \\
\midrule
分解式证明与库配合 & Hilbert、DreamProver 的草图与证明流程 \\
\bottomrule
\end{tabularx}
\end{table}

\subsection{仍然可能是空位的部分}

\begin{enumerate}[leftmargin=1.6em,itemsep=3pt]
  \item \textbf{用实测的工作负载覆盖增益来做库的选择与淘汰。}
        现有工作的判据都是启发式（使用频率、结构相似度、语义聚类）。
        而"覆盖"是并集形式的函数，天然\textbf{单调子模}，因此贪心可给出 $1-1/e$ 的近似保证。
        把"有近似保证的优化"与"启发式排序"放在同一工作负载上对打，是一个干净的、
        可证伪的对比实验。本次 arXiv 检索（主题词含 submodular 与 data selection）
        \textbf{未发现相关工作}，但需注意 arXiv 检索只覆盖元数据，不能视为穷尽。
  \item \textbf{需求驱动：从证明轨迹中挖掘跨目标反复出现的子目标，用它决定"造什么引理"。}
        SGS 只是条件化在未解目标上，并不做需求挖掘与跨目标过滤；
        线 B 的工作基本是"从已产生的中间定理里事后筛选"。这个"事前定方向"的角度未见对应工作。
  \item \textbf{用检索替代容量压缩。}
        DreamProver 因整库注入而必须把库限制在 100 条以内；
        若库改为按需检索，容量约束就从"能否装进上下文"变为"检索精度"，
        库可以长得更大。这是一个尚未被检验的设计选择。
  \item \textbf{形式化判据作为准入与排序。}
        MathlibLemma 有"证明旁路筛查"，DreamProver 有结构相似度阈值，
        但没有人用一套"硬门乘软分"式的可计算判据同时承担准入、排序与淘汰。
\end{enumerate}

\subsection{必须正面处理的两条威胁}

\begin{enumerate}[leftmargin=1.6em,itemsep=3pt]
  \item \textbf{LeanPremise 的"动态适应用户上下文"与我们的库在功能上高度重叠。}
        它是"从用户环境动态抽取签名以支持训练外库与本地引理"，
        我们做的是"积累自己的引理并喂给证明器"。定位时必须给出清晰区别，
        否则容易被判定为重复。
  \item \textbf{"按有用性筛引理"有 2014 年的先例}（HOL Light 与 Flyspeck 的百万级推理图）。
        我们若只把"有用性"换一个名字，不构成新意；
        差异必须在"\textbf{面向工作负载的覆盖最大化 + 近似保证}"这个层面。
\end{enumerate}

\subsection{对当前项目方案的三条修改建议}

\begin{enumerate}[leftmargin=1.6em,itemsep=3pt]
  \item \textbf{检索层的优先级提到最高}，而且从一开始就采用已公开的检索工具
        （如 LeanExplore 这类可自托管或带 API 的引擎），不要自己从零训检索器。
        理由：这条路已经成熟，自研的期望收益低。
  \item \textbf{主基准从 miniF2F 换成有余量的基准}，
        并在报告里按领域分桶。理由：miniF2F 已被做到 99.2\%。
  \item \textbf{把"库的选择层"当作唯一的主攻方向}，
        用"实测覆盖 + 子模保证"与"启发式排序"做正面对照；
        库的其余环节（去重、淘汰、准入）直接沿用已有方案，不重复造。
\end{enumerate}

\appendix
\section{本次检索并入藏的论文清单}

以下 22 篇已下载到文献库目录
\lp{ATP_Papers/08_Library_Premise_Selection} 下，编号接续原有清单（101 起）。

\begin{table}[htbp]
\centering
\footnotesize
\begin{tabularx}{\textwidth}{@{}p{1.0cm}p{2.4cm}Y@{}}
\toprule
\textbf{编号} & \textbf{arXiv} & \textbf{标题} \\
\midrule
101 & 2604.26311 & DreamProver: Evolving Transferable Lemma Libraries via a Wake-Sleep Theorem-Proving Agent \\
\midrule
102 & 2605.13137 & LeanSearch v2: Global Premise Retrieval for Lean 4 Theorem Proving \\
\midrule
103 & 2506.07477 & Premise Selection for a Lean Hammer \\
\midrule
104 & 2506.11085 & LeanExplore: A search engine for Lean 4 declarations \\
\midrule
105 & 2510.15940 & Lean Finder: Semantic Search for Mathlib That Understands User Intents \\
\midrule
106 & 2605.14549 & CSLibPremiseBench: Structure-Guided Premise Retrieval and Label Robustness for Lean 4 \\
\midrule
107 & 2510.23637 & Combining Textual and Structural Information for Premise Selection in Lean \\
\midrule
108 & 2602.05216 & Semantic Search over 9 Million Mathematical Theorems \\
\midrule
109 & 2602.02561 & MathlibLemma: Folklore Lemma Generation and Benchmark for Formal Mathematics \\
\midrule
110 & 2605.17283 & OProver: A Unified Framework for Agentic Formal Theorem Proving \\
\midrule
111 & 2509.22819 & Hilbert: Recursively Building Formal Proofs with Informal Reasoning \\
\midrule
112 & 2304.00994 & Machine-Learned Premise Selection for Lean \\
\midrule
113 & 2303.15642 & Graph Sequence Learning for Premise Selection \\
\midrule
114 & 2501.15797 & LemmaHead: RAG Assisted Proof Generation Using Large Language Models \\
\midrule
115 & 2502.17925 & LeanProgress: Guiding Search for Neural Theorem Proving via Proof Progress Prediction \\
\midrule
116 & 2310.04353 & An In-Context Learning Agent for Formal Theorem-Proving \\
\midrule
117 & 2604.19558 & On Reasoning-Centric LLM-based Automated Theorem Proving \\
\midrule
118 & 2601.14027 & Numina-Lean-Agent: An Open and General Agentic Reasoning System for Formal Mathematics \\
\midrule
119 & 1402.3578 & Learning-assisted Theorem Proving with Millions of Lemmas \\
\midrule
120 & 2505.20613 & REAL-Prover: Retrieval Augmented Lean Prover for Mathematical Reasoning \\
\midrule
121 & 2607.27259 & CircuitProver: Agentic Lean 4 Theorem Proving with Reusable Circuit Proof Library \\
\midrule
122 & 2510.10815 & DRIFT: Decompose, Retrieve, Illustrate, then Formalize Theorems \\
\bottomrule
\end{tabularx}
\end{table}

\textbf{注}：另有若干重要工作位于本地文献库的既有分类中，本次一并引用，包括：

\begin{itemize}[leftmargin=1.6em,itemsep=1pt]
  \item DeepMath，\lp{arXiv:1606.04442}
  \item Isabelle ENIGMA，\lp{arXiv:2205.01981}
  \item Magnushammer，\lp{arXiv:2303.04488}
  \item LeanDojo 与 ReProver，\lp{arXiv:2306.15626}
  \item LEGO-Prover，\lp{arXiv:2310.00656}
  \item MINIMO，\lp{arXiv:2407.00695}
  \item SGS，\lp{arXiv:2604.20209}
\end{itemize}

\end{document}
